\documentclass[11pt]{article}
\usepackage[T1]{fontenc}
\usepackage[utf8]{inputenc}
\usepackage[spanish,es-tabla]{babel}
\usepackage{geometry}
\usepackage{longtable}
\usepackage{booktabs}
\usepackage{tabularx}
\usepackage{array}
\usepackage{enumitem}
\usepackage{hyperref}
\usepackage{xcolor}
\usepackage{fancyhdr}
\geometry{margin=2.2cm}
\hypersetup{hidelinks}
\pagestyle{fancy}
\fancyhf{}
\fancyhead[L]{Guía Técnica del Sistema}
\fancyhead[R]{Procesamiento Distribuido de Imágenes}
\fancyfoot[C]{\thepage}
\setlength{\headheight}{14pt}
\setlist[itemize]{leftmargin=*, itemsep=0.25em, topsep=0.35em}
\setlist[enumerate]{leftmargin=*, itemsep=0.3em, topsep=0.35em}
\renewcommand{\arraystretch}{1.2}

\title{\textbf{Guía Técnica del Sistema}\\Documentación Parte por Parte}
\author{Uso interno del proyecto}
\date{\today}

\begin{document}
\maketitle
\tableofcontents
\newpage

\section{Objetivo de esta guía}
Esta guía explica el sistema completo por partes, con un enfoque práctico. Para cada componente se documenta:
\begin{itemize}
\item qué hace;
\item en qué archivo vive;
\item qué recibe;
\item qué entrega;
\item cómo se conecta con los demás.
\end{itemize}

La intención es que esta guía sirva para estudiar el proyecto, exponerlo o retomarlo después sin tener que reconstruir todo el contexto mental a partir del código.

\section{Vista general del sistema}
El sistema completo está compuesto por los siguientes bloques:
\begin{itemize}
\item \textbf{Cliente gRPC}: envía solicitudes de negocio al sistema.
\item \textbf{HAProxy}: publica un único endpoint para el clúster coordinador.
\item \textbf{Dos coordinadores}: uno activo y uno pasivo, con failover real.
\item \textbf{Redis}: decide qué coordinador es el líder.
\item \textbf{Tres workers}: ejecutan el procesamiento de imágenes.
\item \textbf{MinIO}: almacena estado compartido, resultados y pendientes.
\item \textbf{Backends concretos}: OCR con Tesseract e inferencia con un modelo ligero.
\end{itemize}

\subsection{Resumen del flujo extremo a extremo}
\begin{enumerate}
\item El cliente se conecta a \texttt{127.0.0.1:50052}.
\item HAProxy redirige la conexión al coordinador que esté marcado como \texttt{ready}.
\item El coordinador crea o reutiliza una solicitud lógica global.
\item El coordinador selecciona el worker más conveniente.
\item El worker acepta la tarea, la encola localmente y la ejecuta.
\item El worker reporta progreso, heartbeat y resultado al coordinador.
\item El coordinador consolida el estado y responde al cliente.
\end{enumerate}

\section{Mapa rápido de archivos}
\begin{longtable}{>{\raggedright\arraybackslash}p{3.2cm} >{\raggedright\arraybackslash}p{4.1cm} >{\raggedright\arraybackslash}p{6.8cm}}
\toprule
\textbf{Parte} & \textbf{Archivo principal} & \textbf{Responsabilidad} \\
\midrule
Coordinador & \texttt{coordinator/runtime.py} & Cola global, deduplicación, ownership, despacho y callbacks \\
Liderazgo & \texttt{coordinator/leadership.py} & Elección del líder con Redis \\
Salud del coordinador & \texttt{coordinator/health.py} & \texttt{/livez} y \texttt{/readyz} del coordinador \\
Persistencia coordinador & \texttt{coordinator/state\_store.py} & Estado durable de la cola y de resultados \\
Worker & \texttt{worker/core/node.py} & Orquestación interna del nodo \\
API de negocio del worker & \texttt{worker/core/business\_api.py} & Traducción entre contrato de negocio y tareas internas \\
Scheduler & \texttt{worker/scheduler/*.py} & Prioridad, score y costo estimado \\
Ejecución & \texttt{worker/execution/*.py} & Procesamiento real, recursos y cancelación \\
gRPC del worker & \texttt{worker/grpc/*.py} & Servidores, reporter y mapeos protobuf \\
Persistencia worker & \texttt{worker/core/state\_store.py} & Pendientes, resultados y spool durable \\
Storage abstracto & \texttt{worker/core/storage.py} & Local, \texttt{file://}, \texttt{http(s)} y \texttt{s3://} \\
Telemetría & \texttt{worker/telemetry/*.py} & Logs, métricas, tracing y salud \\
OCR & \texttt{scripts/ocr\_backend.py} & OCR real con Tesseract \\
Inferencia & \texttt{scripts/inference\_backend.py} & Inferencia real con modelo centroidal \\
Seguridad & \texttt{scripts/generate\_dev\_security\_assets.py} & Generación de certificados y secretos de desarrollo \\
Despliegue & \texttt{docker-compose.yml} & Redis, MinIO, coordinadores, HAProxy y workers \\
Balanceador & \texttt{deploy/haproxy.cfg} & Endpoint único y enrutamiento al líder \\
\bottomrule
\end{longtable}

\section{Capa de entrada y despliegue}

\subsection{HAProxy}
\textbf{Archivo principal}: \texttt{deploy/haproxy.cfg}

\textbf{Qué hace}
\begin{itemize}
\item Publica el endpoint único del coordinador en \texttt{127.0.0.1:50052}.
\item Revisa \texttt{/readyz} de cada coordinador.
\item Solo enruta tráfico gRPC al coordinador que sea líder.
\end{itemize}

\textbf{Qué recibe}
\begin{itemize}
\item conexiones gRPC de clientes externos;
\item healthchecks HTTP internos sobre el puerto de salud del coordinador.
\end{itemize}

\textbf{Qué entrega}
\begin{itemize}
\item una ruta de entrada estable al clúster coordinador;
\item failover transparente si cambia el líder.
\end{itemize}

\textbf{Cómo se conecta}
\begin{itemize}
\item escucha delante de \texttt{coordinator1} y \texttt{coordinator2};
\item usa \texttt{/readyz} para determinar cuál coordinador está activo.
\end{itemize}

\subsection{Docker Compose}
\textbf{Archivo principal}: \texttt{docker-compose.yml}

\textbf{Qué hace}
\begin{itemize}
\item define el stack completo del sistema;
\item levanta Redis, MinIO, dos coordinadores, HAProxy y tres workers;
\item inyecta certificados, secretos y variables de entorno.
\end{itemize}

\textbf{Qué recibe}
\begin{itemize}
\item imágenes Docker;
\item secretos montados desde \texttt{.secrets/};
\item configuración del proyecto.
\end{itemize}

\textbf{Qué entrega}
\begin{itemize}
\item un entorno reproducible;
\item puertos públicos para gRPC, salud, métricas y consola MinIO.
\end{itemize}

\textbf{Cómo se conecta}
\begin{itemize}
\item Redis se usa para liderazgo;
\item MinIO se usa para storage compartido;
\item workers y coordinadores se comunican por mTLS.
\end{itemize}

\section{Coordinador}

\subsection{Visión general}
El coordinador es el plano de control global. No procesa imágenes; decide a qué worker mandar cada solicitud y mantiene la visión agregada del clúster.

\subsection{Configuración}
\textbf{Archivo principal}: \texttt{coordinator/config.py}

\textbf{Qué hace}
\begin{itemize}
\item carga variables de entorno del coordinador;
\item define puertos, Redis, storage, workers conocidos, TLS y timeouts.
\end{itemize}

\textbf{Entradas}
\begin{itemize}
\item variables \texttt{COORDINATOR\_*}.
\end{itemize}

\textbf{Salidas}
\begin{itemize}
\item un objeto de configuración tipado para el runtime.
\end{itemize}

\subsection{Arranque}
\textbf{Archivo principal}: \texttt{coordinator/main.py}

\textbf{Qué hace}
\begin{itemize}
\item crea el runtime;
\item registra el servicio gRPC de negocio;
\item registra el servicio de callbacks del coordinador;
\item arranca el servidor HTTP de salud.
\end{itemize}

\textbf{Entradas}
\begin{itemize}
\item configuración ya cargada;
\item credenciales TLS;
\item lista de workers.
\end{itemize}

\textbf{Salidas}
\begin{itemize}
\item proceso coordinador en ejecución.
\end{itemize}

\subsection{Runtime principal}
\textbf{Archivo principal}: \texttt{coordinator/runtime.py}

\textbf{Qué hace}
\begin{itemize}
\item recibe solicitudes de negocio;
\item calcula una clave global por request;
\item deduplica si el trabajo ya existe o si el resultado ya fue generado;
\item mantiene la cola global;
\item asigna ownership temporal por solicitud;
\item selecciona el worker más conveniente;
\item recibe progreso, resultados y heartbeat de los workers;
\item guarda resultados y estado durable.
\end{itemize}

\textbf{Entradas}
\begin{itemize}
\item \texttt{ProcessRequest} del cliente;
\item callbacks \texttt{ReportProgress}, \texttt{ReportResult}, \texttt{Heartbeat};
\item estados de workers obtenidos por polling;
\item liderazgo desde Redis.
\end{itemize}

\textbf{Salidas}
\begin{itemize}
\item respuesta al cliente;
\item despacho al worker por gRPC;
\item persistencia del estado global.
\end{itemize}

\textbf{Cómo se conecta}
\begin{itemize}
\item usa \texttt{ImageNodeService} para hablar con clientes;
\item usa \texttt{WorkerControlService} para hablar con workers;
\item usa \texttt{CoordinatorCallbackService} para recibir reportes;
\item usa Redis para liderazgo;
\item usa MinIO para persistencia compartida.
\end{itemize}

\subsection{Elección de líder}
\textbf{Archivo principal}: \texttt{coordinator/leadership.py}

\textbf{Qué hace}
\begin{itemize}
\item adquiere un lock distribuido en Redis;
\item renueva el lock mientras la instancia siga viva;
\item libera el liderazgo al cerrar;
\item permite que otra instancia asuma el rol si la líder cae.
\end{itemize}

\textbf{Entradas}
\begin{itemize}
\item URL de Redis;
\item \texttt{cluster\_id};
\item \texttt{instance\_id};
\item TTL del lock.
\end{itemize}

\textbf{Salidas}
\begin{itemize}
\item bandera \texttt{is\_leader};
\item identificador del holder actual del lock.
\end{itemize}

\subsection{Salud del coordinador}
\textbf{Archivo principal}: \texttt{coordinator/health.py}

\textbf{Qué hace}
\begin{itemize}
\item expone \texttt{/livez};
\item expone \texttt{/readyz}.
\end{itemize}

\textbf{Regla importante}
\begin{itemize}
\item \texttt{/livez}: proceso vivo;
\item \texttt{/readyz}: coordinador vivo \textbf{y} líder.
\end{itemize}

Esa diferencia es la que permite que Docker considere sanos a ambos coordinadores, mientras HAProxy solo envía tráfico al líder.

\subsection{Persistencia del coordinador}
\textbf{Archivo principal}: \texttt{coordinator/state\_store.py}

\textbf{Qué hace}
\begin{itemize}
\item guarda solicitudes pendientes;
\item guarda resultados ya procesados;
\item restaura el estado al reiniciar.
\end{itemize}

\textbf{Entradas}
\begin{itemize}
\item jobs serializados;
\item records procesados;
\item storage compartido.
\end{itemize}

\textbf{Salidas}
\begin{itemize}
\item cola global durable;
\item cache durable de resultados.
\end{itemize}

\section{Worker}

\subsection{Visión general}
Cada worker es un ejecutor especializado. Su responsabilidad es aceptar tareas, administrarlas en una cola priorizada local, decidir cuándo puede correrlas y ejecutar el pipeline de imagen.

\subsection{Configuración del worker}
\textbf{Archivo principal}: \texttt{worker/config.py}

\textbf{Qué hace}
\begin{itemize}
\item carga puertos, límites, storage, TLS, tracing, OCR e inferencia;
\item decide si el worker exige storage compartido.
\end{itemize}

\subsection{Orquestador interno}
\textbf{Archivo principal}: \texttt{worker/core/node.py}

\textbf{Qué hace}
\begin{itemize}
\item mantiene la cola local;
\item lanza el scheduler;
\item maneja admisión;
\item registra tareas activas y completadas;
\item procesa cancelaciones;
\item coordina retries;
\item actualiza salud, métricas y reporter;
\item restaura pendientes y resultados al reiniciar.
\end{itemize}

\textbf{Entradas}
\begin{itemize}
\item tareas nuevas desde gRPC;
\item snapshots de recursos;
\item resultados del execution manager;
\item eventos de cancelación y shutdown.
\end{itemize}

\textbf{Salidas}
\begin{itemize}
\item decisiones de aceptación o rechazo;
\item progreso y resultados hacia el coordinador;
\item persistencia de estado local o compartido.
\end{itemize}

\subsection{API de negocio del worker}
\textbf{Archivo principal}: \texttt{worker/core/business\_api.py}

\textbf{Qué hace}
\begin{itemize}
\item traduce solicitudes del contrato \texttt{ImageNodeService} a tareas internas;
\item construye \texttt{Task} a partir de filtros y payload;
\item espera el resultado y lo devuelve como bytes o como path.
\end{itemize}

\textbf{Entradas}
\begin{itemize}
\item \texttt{ProcessRequest};
\item \texttt{filters};
\item bytes de imagen.
\end{itemize}

\textbf{Salidas}
\begin{itemize}
\item \texttt{ExecutionResultRecord};
\item bytes o path final para la respuesta gRPC.
\end{itemize}

\subsection{Parser de filtros}
\textbf{Archivo principal}: \texttt{worker/core/filter\_parser.py}

\textbf{Qué hace}
\begin{itemize}
\item interpreta filtros de texto;
\item convierte cada filtro en una \texttt{TransformationSpec};
\item detecta el formato de salida.
\end{itemize}

\textbf{Ejemplos}
\begin{itemize}
\item \texttt{resize:640x480}
\item \texttt{rotate:90}
\item \texttt{watermark:Demo|10|10|white}
\item \texttt{format:webp}
\item \texttt{ocr}
\item \texttt{inference}
\end{itemize}

\section{Scheduler local}

\subsection{Cola priorizada}
\textbf{Archivo principal}: \texttt{worker/scheduler/priority\_queue.py}

\textbf{Qué hace}
\begin{itemize}
\item encola tareas;
\item permite sacarlas en orden de prioridad efectiva;
\item soporta remoción y espera asíncrona.
\end{itemize}

\subsection{Scoring}
\textbf{Archivo principal}: \texttt{worker/scheduler/scoring.py}

\textbf{Qué hace}
\begin{itemize}
\item calcula el score de cada tarea;
\item mezcla prioridad, aging, deadline, costo y tamaño.
\end{itemize}

\textbf{Idea del score}
\begin{itemize}
\item las tareas urgentes suben;
\item las viejas ganan aging;
\item las caras o grandes bajan;
\item los retries penalizan;
\item se evita starvation.
\end{itemize}

\subsection{Costo estimado}
\textbf{Archivo principal}: \texttt{worker/scheduler/cost\_model.py}

\textbf{Qué hace}
\begin{itemize}
\item estima el tiempo de servicio;
\item aprende de ejecuciones pasadas usando promedio exponencial.
\end{itemize}

\section{Ejecución y procesamiento de imágenes}

\subsection{Gestión de recursos}
\textbf{Archivo principal}: \texttt{worker/execution/resource\_manager.py}

\textbf{Qué hace}
\begin{itemize}
\item estima requerimientos de una tarea;
\item toma snapshots de CPU, RAM y ocupación;
\item decide si una tarea puede correr ahora.
\end{itemize}

\textbf{Entradas}
\begin{itemize}
\item tarea;
\item longitud de cola;
\item número de tareas activas;
\item métricas del host.
\end{itemize}

\textbf{Salidas}
\begin{itemize}
\item \texttt{ResourceRequirements};
\item \texttt{ResourceSnapshot};
\item decisión de admisión.
\end{itemize}

\subsection{Execution manager}
\textbf{Archivo principal}: \texttt{worker/execution/execution\_manager.py}

\textbf{Qué hace}
\begin{itemize}
\item arranca la ejecución real;
\item limita concurrencia con semáforos;
\item decide entre thread y proceso dedicado;
\item maneja cancelación;
\item observa tareas activas;
\item limpia salidas parciales.
\end{itemize}

\textbf{Entradas}
\begin{itemize}
\item \texttt{Task};
\item requerimientos estimados;
\item tokens de cancelación;
\item configuración de OCR e inferencia.
\end{itemize}

\textbf{Salidas}
\begin{itemize}
\item \texttt{ExecutionResultRecord};
\item estados de cancelación o error.
\end{itemize}

\subsection{Procesador de imágenes}
\textbf{Archivo principal}: \texttt{worker/execution/image\_processor.py}

\textbf{Qué hace}
\begin{itemize}
\item carga la imagen de entrada;
\item aplica transformaciones;
\item ejecuta OCR o inferencia si la tarea lo pide;
\item serializa el formato final;
\item persiste el resultado.
\end{itemize}

\textbf{Operaciones integradas}
\begin{itemize}
\item grayscale;
\item resize;
\item crop;
\item rotate;
\item flip;
\item blur;
\item sharpen;
\item brightness/contrast;
\item watermark/text;
\item conversión de formato.
\end{itemize}

\textbf{Formatos soportados}
\begin{itemize}
\item JPG
\item PNG
\item TIF
\item WEBP
\item BMP
\item GIF estático
\item ICO
\end{itemize}

\textbf{Punto importante}
Las operaciones costosas integradas fueron implementadas de manera cooperativa, para que la cancelación normal no dependa de matar procesos en la mayoría de los casos.

\section{Persistencia y storage}

\subsection{Storage abstracto}
\textbf{Archivo principal}: \texttt{worker/core/storage.py}

\textbf{Qué hace}
\begin{itemize}
\item unifica acceso a distintos backends;
\item lee, escribe, lista y borra URIs;
\item soporta local, \texttt{file://}, \texttt{http(s)} y \texttt{s3://}.
\end{itemize}

\textbf{Por qué existe}
Evita que el resto del sistema tenga que saber si el resultado quedó en disco local o en MinIO.

\subsection{Persistencia del worker}
\textbf{Archivo principal}: \texttt{worker/core/state\_store.py}

\textbf{Qué hace}
\begin{itemize}
\item guarda resultados terminales;
\item guarda tareas pendientes;
\item guarda eventos de progreso y resultado mientras esperan \texttt{ack}.
\end{itemize}

\textbf{Subcomponentes}
\begin{itemize}
\item \texttt{StateStore}: cache durable de completados;
\item \texttt{PendingTaskStore}: tareas en cola o retry;
\item \texttt{EventSpoolStore}: cola durable de eventos hacia el coordinador.
\end{itemize}

\section{gRPC}

\subsection{Contrato de negocio}
\textbf{Archivo principal}: \texttt{proto/imagenode.proto}

\textbf{Qué hace}
\begin{itemize}
\item define la interfaz que usa el cliente para pedir procesamiento;
\item expone modos unary, streaming y batch.
\end{itemize}

\textbf{Métodos principales}
\begin{itemize}
\item \texttt{ProcessToPath}
\item \texttt{ProcessToData}
\item \texttt{UploadLargeImage}
\item \texttt{StreamBatchProcess}
\item \texttt{ProcessBatch}
\item \texttt{HealthCheck}
\item \texttt{GetMetrics}
\end{itemize}

\subsection{Contrato interno distribuido}
\textbf{Archivo principal}: \texttt{proto/worker\_node.proto}

\textbf{Qué hace}
\begin{itemize}
\item define el control entre coordinador y workers;
\item define callbacks del worker al coordinador.
\end{itemize}

\textbf{Servicios}
\begin{itemize}
\item \texttt{WorkerControlService}
\item \texttt{CoordinatorCallbackService}
\end{itemize}

\subsection{Servidor de control del worker}
\textbf{Archivo principal}: \texttt{worker/grpc/servicer.py}

\textbf{Qué hace}
\begin{itemize}
\item implementa \texttt{SubmitTask};
\item implementa \texttt{GetNodeStatus};
\item implementa \texttt{CancelTask};
\item implementa \texttt{DrainNode};
\item implementa \texttt{ShutdownNode}.
\end{itemize}

\subsection{Servidor de negocio del worker}
\textbf{Archivo principal}: \texttt{worker/grpc/business\_servicer.py}

\textbf{Qué hace}
\begin{itemize}
\item expone \texttt{ImageNodeService};
\item usa internamente \texttt{business\_api.py};
\item no crea otra lógica paralela: reutiliza el mismo motor del worker.
\end{itemize}

\subsection{Reporter}
\textbf{Archivo principal}: \texttt{worker/grpc/reporter.py}

\textbf{Qué hace}
\begin{itemize}
\item mantiene el canal hacia el coordinador;
\item envía \texttt{ProgressEvent};
\item envía \texttt{ExecutionResult};
\item envía \texttt{Heartbeat};
\item reintenta con backoff si el coordinador no responde.
\end{itemize}

\textbf{Punto importante}
Si el coordinador no está disponible, los eventos no se pierden: se quedan en el spool durable y se reenvían después.

\subsection{Seguridad gRPC}
\textbf{Archivo principal}: \texttt{worker/grpc/security.py}

\textbf{Qué hace}
\begin{itemize}
\item construye credenciales TLS del servidor;
\item construye credenciales TLS del cliente;
\item soporta mTLS en ambos sentidos.
\end{itemize}

\section{Telemetría}

\subsection{Logs}
\textbf{Archivo principal}: \texttt{worker/telemetry/logging.py}

\textbf{Qué hace}
\begin{itemize}
\item emite logs JSON;
\item incluye contexto útil como tarea, nodo y trazas.
\end{itemize}

\subsection{Métricas}
\textbf{Archivo principal}: \texttt{worker/telemetry/metrics.py}

\textbf{Qué hace}
\begin{itemize}
\item expone métricas Prometheus;
\item mide cola, latencias, retries, tareas activas y conectividad.
\end{itemize}

\subsection{Salud}
\textbf{Archivo principal}: \texttt{worker/telemetry/health.py}

\textbf{Qué hace}
\begin{itemize}
\item expone \texttt{/livez};
\item expone \texttt{/readyz}.
\end{itemize}

\textbf{Diferencia}
\begin{itemize}
\item \texttt{/livez}: proceso vivo;
\item \texttt{/readyz}: worker realmente apto para recibir más carga.
\end{itemize}

\subsection{Tracing}
\textbf{Archivo principal}: \texttt{worker/telemetry/tracing.py}

\textbf{Qué hace}
\begin{itemize}
\item propaga contexto de trazas;
\item integra OpenTelemetry cuando está habilitado.
\end{itemize}

\section{Backends concretos}

\subsection{OCR}
\textbf{Archivo principal}: \texttt{scripts/ocr\_backend.py}

\textbf{Qué hace}
\begin{itemize}
\item lee una imagen temporal;
\item invoca Tesseract vía \texttt{pytesseract};
\item devuelve texto reconocido en JSON.
\end{itemize}

\textbf{Entradas}
\begin{itemize}
\item \texttt{WORKER\_INPUT\_FILE}
\item \texttt{WORKER\_OUTPUT\_JSON}
\item parámetros opcionales como idioma y \texttt{psm}
\end{itemize}

\textbf{Salidas}
\begin{itemize}
\item texto OCR;
\item metadata del motor usado.
\end{itemize}

\subsection{Inferencia}
\textbf{Archivo principal}: \texttt{scripts/inference\_backend.py}

\textbf{Qué hace}
\begin{itemize}
\item extrae features visuales de la imagen;
\item compara contra centroides;
\item devuelve etiquetas y score.
\end{itemize}

\textbf{Modelo}
\begin{itemize}
\item el artefacto vive en \texttt{worker/backends/inference\_model.json};
\item es un modelo ligero y concreto incluido en el repo.
\end{itemize}

\section{Seguridad operativa}

\subsection{Generación de PKI y secretos}
\textbf{Archivo principal}: \texttt{scripts/generate\_dev\_security\_assets.py}

\textbf{Qué hace}
\begin{itemize}
\item crea una CA de desarrollo;
\item genera certificados para coordinador, workers y cliente demo;
\item genera secretos de acceso para MinIO.
\end{itemize}

\textbf{Salidas}
\begin{itemize}
\item \texttt{.secrets/pki/*.crt}
\item \texttt{.secrets/pki/*.key}
\item \texttt{.secrets/minio\_access\_key.txt}
\item \texttt{.secrets/minio\_secret\_key.txt}
\end{itemize}

\section{Scripts y herramientas de operación}

\subsection{run.ps1}
\textbf{Archivo principal}: \texttt{run.ps1}

\textbf{Qué hace}
\begin{itemize}
\item instala dependencias;
\item regenera protos;
\item levanta el stack;
\item ejecuta la demo;
\item baja el stack;
\item corre tests.
\end{itemize}

\textbf{Comandos más importantes}
\begin{itemize}
\item \texttt{.\textbackslash run.ps1 stack}
\item \texttt{.\textbackslash run.ps1 demo}
\item \texttt{.\textbackslash run.ps1 down}
\item \texttt{.\textbackslash run.ps1 test}
\end{itemize}

\subsection{Demo extremo a extremo}
\textbf{Archivo principal}: \texttt{scripts/demo\_end\_to\_end.py}

\textbf{Qué hace}
\begin{itemize}
\item construye una imagen de prueba;
\item llama al sistema por TLS;
\item valida \texttt{HealthCheck};
\item hace una solicitud simple y una por batch;
\item guarda artefactos en \texttt{docs/demo/}.
\end{itemize}

\section{Cómo seguir una solicitud real}

\subsection{Camino normal}
\begin{enumerate}
\item El cliente invoca \texttt{ImageNodeService}.
\item El coordinador recibe la request en \texttt{coordinator/runtime.py}.
\item Se calcula un \texttt{request\_key}.
\item Si la tarea ya existe o el resultado ya está, se reutiliza.
\item Si es nueva, entra a la cola global.
\item El coordinador elige un worker.
\item El worker recibe la \texttt{Task} por \texttt{SubmitTask}.
\item El worker la encola localmente.
\item El scheduler local la saca cuando hay recursos.
\item El execution manager corre el pipeline.
\item El worker persiste resultado y lo reporta.
\item El coordinador actualiza su estado y responde.
\end{enumerate}

\subsection{Camino de failover}
\begin{enumerate}
\item El coordinador líder cae.
\item Redis deja expirar o libera el lock.
\item El follower adquiere el liderazgo.
\item Su \texttt{/readyz} pasa a \texttt{200}.
\item HAProxy empieza a enviarle tráfico.
\item Los clientes siguen usando el mismo endpoint público.
\end{enumerate}

\section{Qué parte decide qué}
\begin{longtable}{>{\raggedright\arraybackslash}p{3.6cm} >{\raggedright\arraybackslash}p{9.8cm}}
\toprule
\textbf{Componente} & \textbf{Decisión principal} \\
\midrule
HAProxy & A qué coordinador enviar tráfico de entrada \\
Redis + leadership & Qué coordinador es el líder \\
CoordinatorRuntime & Qué solicitud existe, a qué worker despacharla y qué resultado devolver \\
WorkerNode & Si acepta, encola, reintenta, cancela o drena una tarea \\
TaskScorer & Qué prioridad efectiva tiene la tarea \\
ResourceManager & Si hay recursos suficientes para correrla \\
ExecutionManager & Cómo ejecutarla y cómo cancelarla \\
ImageProcessor & Qué transformaciones aplicar y cómo producir el output final \\
State stores & Qué estado debe sobrevivir a reinicios \\
\bottomrule
\end{longtable}

\section{Puntos clave para explicar el proyecto}
Si necesitas resumir el sistema oralmente, estas son las ideas más importantes:
\begin{itemize}
\item el coordinador decide y el worker ejecuta;
\item Redis elige al líder del coordinador;
\item HAProxy da un único endpoint al cliente;
\item MinIO comparte estado y resultados entre nodos;
\item el worker tiene cola local con prioridad, backpressure y persistencia;
\item el sistema usa mTLS y soporta OCR e inferencia reales.
\end{itemize}

\section{Conclusión}
El proyecto quedó dividido de manera clara entre plano de control y plano de ejecución. El coordinador se encarga de la coordinación global, el ownership y el despacho; cada worker se encarga de priorizar, admitir, ejecutar, persistir y reportar. Redis y HAProxy permiten failover real del coordinador, mientras MinIO resuelve la durabilidad compartida. Con esta separación, el sistema es entendible, demostrable y extensible sin rehacer su arquitectura base.

\end{document}
